Observations of the large-scale structure of the Universe, combining weak gravitational lensing with galaxy clustering, provide a powerful probe of cosmology \citep{1972gcpa.book.....W, 1980lssu.book.....P, 1990ApJ...349L...1T, 1998ApJ...498...26K}. Over the past decade, wide-field imaging surveys such as the Kilo-Degree Survey \citep[KiDS;][]{2019A&A...625A...2K, 2024A&A...686A.170W}, the Dark Energy Survey \citep[DES;][]{2025MNRAS.543.4156Y, 2026ApJS..282...62B}, and the Hyper Suprime-Cam Survey \citep[HSC;][]{2022PASJ...74..247A, 2022PASJ...74..421L} have delivered increasingly precise measurements of the auto- and cross-correlations of galaxy shapes and positions \citep{2022PhRvD.105b3514A, 2022PhRvD.105b3515S, 2023PhRvD.108l3519D, 2023PhRvD.108l3518L, 2025A&A...702A.169S, 2025A&A...703A.158W, 2026arXiv260114559A, 2026arXiv260210065D}. These two-point functions form the foundation of the now-standard \(3 \times 2 \,\mathrm{pt}\) framework \citep{2021A&A...646A.140H, 2022PhRvD.105b3520A, 2023A&A...675A.189D, 2023PhRvD.108l3521S, 2023PhRvD.108l3517M}, enabling stringent constraints on late-time structure growth and the cosmic expansion history at a level of precision increasingly competitive with early-Universe measurements from the cosmic microwave background \citep[CMB;][]{2020A&A...641A...6P}.

The next generation of Stage-IV facilities---the Vera C.~Rubin Observatory \citep[\textit{Rubin}\footnote{\url{https://rubinobservatory.org}};][]{2009arXiv0912.0201L}, \textit{Euclid}\footnote{\url{https://www.cosmos.esa.int/web/euclid}} \citep{2025A&A...697A...1E}, the Nancy Grace Roman Space Telescope \citep[\textit{Roman}\footnote{\url{https://roman.gsfc.nasa.gov}};][]{2021MNRAS.507.1514E}, and the China Space Station Telescope \citep[CSST\footnote{\url{http://www.bao.ac.cn/csst/}};][]{2019ApJ...883..203G}---will extend sky coverage, depth, and image quality by orders of magnitude. A central ingredient of the resulting analyses is the computation of angular power spectra under the Limber approximation \citep{1953ApJ...117..134L, 2008PhRvD..78l3506L}, which underpins not only the \(3 \times 2 \,\mathrm{pt}\) framework but any inference based on projected anisotropies of the matter and tracer fields \citep{2018ARA&A..56..393M}. These projections must be evaluated repeatedly---often millions of times \citep[e.g.][]{2024OJAp....7E..73P}---across parameter spaces that span cosmology, astrophysical modelling, and survey-level systematics, placing severe demands on theoretical pipelines.

In current analyses, the cost of Limber integrals grows rapidly with multipole resolution, the number of tomographic bins, and the complexity of the nuisance-parameter model. Modern pipelines incorporate detailed baryonic-feedback prescriptions \citep[e.g.][]{2019OJAp....2E...4C, 2020MNRAS.495.4800A, 2021MNRAS.502.1401M, 2021JCAP...12..046G, 2025MNRAS.539.1337S}, scale-dependent galaxy-bias models \citep[e.g.][]{2018PhR...733....1D, 2020MNRAS.492.5754M, 2021MNRAS.505.1422K}, flexible intrinsic-alignment models \citep[e.g.][]{2013MNRAS.431..477J, 2013MNRAS.436..819J, 2015PhR...558....1T, 2019PhRvD.100j3506B, 2021MNRAS.501.2983F, 2024OJAp....7E..45V}, and sophisticated treatments of tomographic redshift distribution uncertainties \citep[e.g.][]{2020JCAP...10..056H, 2021A&A...650A.148S, 2022MNRAS.511.2170C, 2023MNRAS.522.5037R, 2023OJAp....6E..23H, 2025arXiv250600758B}. Each extension enlarges the parameter space, and the forward-modelling step---recomputing angular power spectra at every likelihood evaluation---often dominates the total inference cost.

A range of strategies have been developed to alleviate this bottleneck. Emulators of the non-linear matter power spectrum \citep{2019MNRAS.484.5509E, 2021MNRAS.505.2840E, 2017ApJ...847...50L, 2023MNRAS.520.3443M, 2019ApJ...875...69D, 2023JCAP...07..054D, 2025SCPMA..6889512C}, together with recent factorisation-based decompositions of the matter power spectrum into a small set of cosmology-independent basis functions \citep{2025PhRvL.134s1002B}, accelerate the three-dimensional modelling stage but do not directly emulate angular power spectra, which depend on survey-specific selection functions and tomographic redshift distributions. GPU-enabled cosmology frameworks based on JAX \citep{2021ascl.soft11002B} and Julia \citep{2017SIAMR..59...65B}---including \texttt{JAX-Cosmo} \citep{2023OJAp....6E..15C} and \texttt{LimberJack.jl} \citep{2024OJAp....7E..11R}---accelerate the line-of-sight projection itself, and additionally enable gradient-based samplers that reduce the number of likelihood evaluations required. Complementary approaches reformulate the Limber and beyond-Limber integrals using FFTLog \citep{2000MNRAS.312..257H} expansions or specialised basis functions to handle oscillatory kernels efficiently \citep[e.g.][]{2017A&A...602A..72C, 2017JCAP...11..054A, 2018JCAP...10..047S, 2018PhRvD..97b3504G, 2020JCAP...05..010F, 2023PhRvD.108j3007F, 2025JOSS...10.8618R}; recent implementations such as \texttt{Blast.jl} \citep{2024OJAp....7E.112C} and \(\texttt{Swift}\,C_\ell\) \citep{2025arXiv250522718R} combine these techniques with modern hardware to achieve significant runtime reductions. At the sampler level, frameworks such as \texttt{Cobaya} \citep{2021JCAP...05..057T} and \texttt{CosmoSIS} \citep{2015A&C....12...45Z} exploit the observation that nuisance parameters typically do not affect the three-dimensional power spectrum, enabling a fast--slow decomposition in which the expensive cosmological computation is reused across multiple nuisance-parameter proposals. However, all of these approaches intertwine cosmology-dependent quantities with survey-specific components such as tomographic binning and tomographic redshift distributions, so the Limber projection must still be re-evaluated whenever any nuisance parameter---in particular the tomographic redshift distributions---changes, even when the underlying cosmology is held fixed.

A formulation that cleanly separates these two classes of dependence would enable direct emulation of the cosmological projection coefficients and open the door to analytic marginalisation over nuisance parameters, substantially reducing the burden on high-dimensional sampling \citep{2020JCAP...10..056H, 2023MNRAS.522.5037R, 2023OJAp....6E..23H, 2025A&A...693A.178R, 2026arXiv260424425R}. In this work, we present \texttt{LimberCloud}, a tensorised analytic reformulation of the Limber projection that achieves precisely this separation. By piecewise-linearising the radial kernels and the matter power spectrum on a discrete redshift grid, the standard Limber integrals reduce to compact tensor contractions of the form \(C^{ab}_\ell = \sum_{ij} B^{ij}_\ell \, \phi^{a}_{i} \, \phi^{b}_{j}\), where the coefficients \(B^{ij}_\ell\) depend only on cosmology and the tomographic redshift distribution vectors \(\phi^{a}_{i}\) and \(\phi^{b}_{j}\) encode all survey-specific information. This factorisation is exact in the limit of infinitely fine redshift binning and achieves sub-per-cent agreement with the Core Cosmology Library \citep[\texttt{CCL};][]{2019ApJS..242....2C} for the bin widths adopted in current Stage-IV forecasts. Implemented in JAX, the framework exploits just-in-time compilation, automatic vectorisation, and GPU acceleration; in the large-batch regime relevant for inference, the JAX--GPU implementation outperforms \texttt{CCL} by more than an order of magnitude. The tensorised structure further provides a natural foundation for future end-to-end emulation of the cosmology-dependent coefficients and for analytic marginalisation over nuisance parameters, both of which we leave to follow-up work. The present work focuses on the mathematical formulation, its validation against \texttt{CCL}, and computational benchmarks across CPU and GPU architectures.

The remainder of this paper is organised as follows. Section~\ref{sec:2} reviews the theoretical formalism of angular power spectra under the Limber approximation. Section~\ref{sec:3} describes the validation setup. Section~\ref{sec:4} introduces the linearisation of the kernel functions and matter power spectra. Section~\ref{sec:5} presents the tensorised reformulation of the Limber projection and compares its accuracy with benchmark calculations. Section~\ref{sec:6} evaluates performance on CPU and GPU architectures. Section~\ref{sec:7} summarises the framework, discusses analytic marginalisation prospects, and outlines future directions.
